\documentclass[10pt]{article}
\usepackage[margin=20mm]{geometry}
\usepackage{amsmath,amssymb,mathtools,fontspec,unicode-math,xurl,hyperref}
\IfFontExistsTF{DejaVu Serif}{\setmainfont{DejaVu Serif}}{\setmainfont{Latin Modern Roman}}
\setmathfont{latinmodern-math.otf}
\setlength{\parindent}{0pt}\setlength{\parskip}{6pt}\sloppy
\title{Technical verification of group theoretical characterizations of rationality}
\author{AI-assisted technical review}
\date{}
\begin{document}
\maketitle
Abstract: This report examines the published paper’s hypotheses, selected descent and derived-length arguments, and the final rationality/ruledness reductions. One quantified auxiliary proposition fails for the trivial algebraic group. A second, local boundary problem occurs in the stated branch-polynomial construction in dimension two. Both have explicit, limited repairs, and neither establishes a failure of the principal rationality or ruledness theorems. The supported assessment is Minor. The review does not certify every imported theorem or every proof in the paper; scientific adjudication remains pending.\par
\section*{Paper identity and scope}
The reviewed work is Andriy Regeta, Christian Urech, and Immanuel van Santen, Group theoretical characterizations of rationality, Inventiones mathematicae 241 (2025), 681–716, DOI 10.1007/s00222-025-01348-7 (https://doi.org/10.1007/s00222-025-01348-7) [paper]. The source is the acquired 36-page publisher Version of Record. PDF page \(j\) corresponds to printed page \(680+j\). Statements below concern those exact bytes.\par
The full extracted article was read to map the argument. Detailed checks concerned the characteristic-zero interpolation in Proposition 4.2, the countable-index and dimension argument leading to Theorem 1.3, Proposition 6.7 and its applications, the boundary construction in Corollary 1.4, and the divisibility obstruction in Proposition 6.9. The Rosenlicht construction and the geometric descent machinery were read as dependencies but not independently re-proved in their entirety. In particular, this review does not independently establish the paper’s imported results on algebraic families, regularization, or rational additive-group actions.\par
\section*{Definitions and checks of the principal reductions}
The standing base field in Sections 5–6 is algebraically closed, uncountable, and of characteristic zero (PDF p. 4 and the relevant section openings). A subgroup is not, by convention, required to be nontrivial: the paper explicitly adds that requirement in statements such as Theorem 6.1. Derived length means the least \(i\) for which the \(i\)-th derived subgroup is trivial, so the identity group has length zero.\par
\subsection*{Descent in Proposition 4.2}
On PDF pp. 15–16 (printed pp. 695–696), a rational additive action defined over a finite Galois extension is descended by evaluating its rational comorphism at integer times. Let \(F=k(X)\) and let \(L(X)/F\) be that extension. After multiplying a denominator by its conjugates, one obtains \(q(t)\in F[t]\). If \(p(t)=\sum_{j=0}^{m}p_jt^j\) takes values in \(F\) at \(m+1\) consecutive sufficiently large integers, then the coefficient vector lies in \(F^{m+1}\), because \[\det\big((N+i)^j\big)_{0\le i,j\le m}
   =\prod_{0\le i<j\le m}(j-i)\ne0.\] Characteristic zero is essential to this computation and is assumed. The finite set of denominator zeros can be avoided. Thus the interpolation step itself does not hide a positive-characteristic extrapolation. The converse also uses that an additive action whose time-one element fixes \(k(Y)\) fixes it at all integer times and hence identically, by the same rational-function identity principle. This checks the descent mechanism, conditional on the cited correspondence between rationally integrable derivations and rational additive actions.\par
\subsection*{Derived length and the equality case}
Lemmas 5.5–5.8 and the proof of Theorem 1.3 are on PDF pp. 19–21. The use of uncountability is explicit: an irreducible parameter variety cannot be partitioned into countably many proper closed coset preimages. Countable-index subgroups are consequently dense, allowing the derived-length comparison.\par
Writing \(d=\dim X-\dim Y\) for the generic-fibre dimension, the induction gives \[\operatorname{len}(G)\le \operatorname{len}(H^\circ)+\operatorname{len}(G/H^\circ)
 \le 2d+2\dim Y=2\dim X.\] If equality holds, both summands attain their individual bounds. Conjugation on the self-centralizing commutative algebraic group on the generic fibre then gives \(2d\le 1+d\); since \(d>0\), this forces \(d=1\). The subsequent additive-group descent and triangular reduction use this one-dimensional case. Closure bars that disappear in plain-text extraction were checked against the rendered source, especially on PDF p. 21. No contradiction in this chain was established.\par
\subsection*{The divisibility obstruction}
Proposition 6.9 (PDF p. 31) states that for \(\dim X>0\) some \(f\in k(X)^\times\) has no positive power divisible by every positive integer. For a transcendence basis, put \(F=k(t_1,\ldots,t_n)\) and \(D=[k(X):F]\). Given \(l\ge1\), choose a prime \(p>l\) with \(p\nmid D\). The \(t_1\)-adic valuation rules out a \(p\)-th root of \(t_1^l\) in \(F\); the cited Kummer lemma then gives degree \(p\) for such a root. It cannot lie in an extension of degree \(D\). This verifies the specific obstruction used in Case 2 of Theorem 1.1. The existence of divisible powers in the auxiliary closed affine algebraic subgroup is imported by the paper from its reference [33], and was not independently proved here.\par
\section*{Local findings}
Problem 1 (The trivial group violates Proposition 6.7).  The statement of Proposition 6.7 on PDF p. 29 (printed p. 709) asserts that \(\operatorname{rRepl}_X(G)\) is self-centralizing for a connected commutative affine algebraic subgroup \(G\subseteq\operatorname{Bir}(X)\), without excluding \(G=\{1\}\).\par
\subsection*{Exact counterexample and hypotheses}
Take \(k=\mathbb C\), \(X=\mathbb P^1_{\mathbb C}\), and \(G=\{\mathrm{id}_X\}\). The variety is irreducible and the base field satisfies every standing assumption. The one-point algebraic group is affine, connected, commutative, and is an algebraic subgroup of \(\operatorname{Bir}(X)\). Its inclusion is faithful.\par
By Definition 6.5 on PDF p. 28, a rational replica comes from a \(G\)-invariant rational map \(X\dashrightarrow G\). There is exactly one such map, and it induces the identity. Therefore \[\operatorname{rRepl}_X(G)=\{\mathrm{id}_X\},\qquad
 \operatorname{Cent}_{\operatorname{Bir}(X)}(\operatorname{rRepl}_X(G))=\operatorname{Bir}(X).\] The map \(z\mapsto z+1\) lies in the right-hand group and is not the identity. The asserted equality is consequently false for an input satisfying the written hypotheses. This is an exact group-theoretic counterexample; no numerical or author-code execution is involved.\par
\subsection*{Counterevidence, repair, and impact}
The proof of Claim 3 on PDF p. 30 chooses a morphism to \(G'\) that separates two points of \(Y\). That step cannot work when \(G'=\{1\}\). The counterexample therefore identifies an actual missing hypothesis, rather than a rendering problem. A precise repair is to assume \(G\ne\{1\}\) in Proposition 6.7. In its proof, a nontrivial connected group acting faithfully on an irreducible variety has a positive-dimensional generic orbit; the resulting \(G'\) is nontrivial. Rational maps to its split additive or multiplicative factors can then separate the specified points after shrinking the domain. This supplies the needed nondegeneracy for Claim 3; it is not a new proof of the imported quotient theorem.\par
The applications in Theorem 1.1 use a nontrivial vector group \(U\) containing the nontrivial image of translations in its replica group, and a positive-dimensional torus \(D\). Theorem 1.2 uses \(U=\mathbb G_a\) and \(D=\mathbb G_m\) explicitly (PDF pp. 33–34). The displayed counterexample does not reach these applications. The defect is therefore Minor: a false boundary instance of an auxiliary proposition, with no demonstrated damage to the principal conclusions.\par
Problem 2 (The branch-polynomial choice excludes \(n=2\)).  Corollary 1.4 includes every \(n\ge2\) (PDF p. 3). Its proof on PDF p. 23 (printed p. 703) begins by choosing a homogeneous irreducible polynomial \(F\in k[x_0,\ldots,x_{n-1}]\) of degree \(2n\) with smooth projective zero locus.\par
\subsection*{Boundary check and a valid replacement}
For \(n=2\), this would be an irreducible binary quartic over an algebraically closed field. Every binary homogeneous polynomial of degree four factors into four linear forms, counted with multiplicity: dehomogenize at \(x_0=1\), factor the resulting univariate polynomial, and include any missing degree as a power of \(x_0\). Thus the requested irreducible polynomial does not exist. This prevents the written construction from starting in that dimension; it does not disprove the corollary.\par
A concrete repair for this boundary case is the squarefree quartic \[F(x_0,x_1)=x_1(x_1-x_0)(x_1-2x_0)(x_1-3x_0).\] Its branch divisor in \(\mathbb P^1\) consists of four distinct finite points; \(x_0\) is not a factor, so the original projective hypersurface gains no extra component at infinity. Its dehomogenization is \(f(y)=y(y-1)(y-2)(y-3)\), which is not a square in \(k(y)\), since its valuation at \(y=0\) is one. The double cover \(z^2=f(y)\), completed over \(\mathbb P^1\), is smooth and has \[2g-2=2(-2)+4=0,\qquad g=1.\] It is therefore non-ruled. The matrices defining \(T_f(\mathbb A^1)\) have determinant \(a^2-b^2f\ne0\) for \((a,b)\ne(0,0)\), because \(f\) is not a square. Hence the nonsquare and non-ruled conditions actually used later in the construction survive this replacement.\par
\subsection*{Counterevidence and impact boundary}
The article already identifies \(n=2\) as a previously proved case, citing its reference [20] both in the introduction and surrounding discussion. Reading the new construction as intended only for \(n\ge3\) therefore supplies a second resolution: state that restriction and retain the cited two-dimensional result. That earlier paper was not independently audited here.\par
The supported finding is the local incompatibility of the written construction with its endpoint, not a new counterexample to the existence of nonstandard Borel subgroups. The corrected squarefree construction was checked above; the higher-dimensional construction is not challenged by this argument. Its impact is Minor, limited to presentation of the already known two-dimensional endpoint.\par
\section*{Overall assessment}
The two findings concern different boundaries: the identity group in a general auxiliary proposition, and an impossible irreducible branch polynomial at an endpoint dimension. The checked applications of Proposition 6.7 avoid the first boundary. The paper itself cites prior treatment of the second, and an explicit squarefree replacement retains the required geometry. Neither finding establishes an invalid rationality or ruledness theorem, nor does either warrant a Major assessment.\par
Item | Assessment | Repair verification and remaining state\par
Proposition 6.7 | Minor | Exact counterexample; nontriviality repairs the identified boundary.\par
Corollary 1.4 construction | Minor | Explicit \(n=2\) squarefree replacement checked.\par
Principal reductions | No issue established in checked steps | Imported results and the full proof network not certified.\par
Scientific decision | Pending | No human review or author adoption recorded.\par
\section*{Supporting checks}
> 592100d26bb0f018827356cc47a98c89e7eea9a66151f2fef20e4278df66f507.\par
The source manifest, extracted text, inspected source-page images, complete coverage ledger, findings, and actual build/inspection records are retained with this paper’s verification files. All arguments in this report are analytic. No empirical reproduction or execution of author code is claimed. The source-page images were examined to distinguish actual mathematical notation from text-extraction losses.\par
This is an initial AI technical review in the assigned recovery session, not independent acceptance. The original task identity and recovery identity are preserved in provenance; the source-acquisition token is kept as historical acquisition metadata. No later version or erratum was used to infer correction or author adoption. A dedicated later-version literature audit was not performed.\par
1 A. Regeta, C. Urech, and I. van Santen, Group theoretical characterizations of rationality, Invent. Math. 241 (2025), 681–716. doi:10.1007/s00222-025-01348-7 (https://doi.org/10.1007/s00222-025-01348-7). Exact acquired publisher PDF, 36 pages.\par
\end{document}
